\chapter{Créditos y Asignaturas}

\section{Créditos}

La \acrshort{UDC} sigue el sistema \textbf{\acrfull{ECTS}}.

\begin{curiosityBox}
    El sistema \acrshort{ECTS} forma parte del Plan Bolonia y es utilizado por las universidades europeas para convalidar asignaturas y \textbf{cuantificar el trabajo relativo al estudiante} que trabaja bajo los grados del \acrfull{EEES}.
\end{curiosityBox}

Un \textbf{crédito} suele representar un \textbf{número determinado de horas} de trabajo del estudiante, incluyendo tanto las horas de clase como el tiempo dedicado al estudio y la realización de tareas y proyectos fuera del aula. La correspondencia entre asignaturas y créditos varía en función de la carga de trabajo que se espera del estudiante.

En general, los planes de estudio establecen un \textbf{número mínimo de créditos} que el estudiante debe completar para obtener su título.

\begin{infoBox}
    Los tres grados de la \acrshort{FIC} (\href{\linkPortalGEI}{\acrshort{GEI}}, \href{\linkPortalGCED}{\acrshort{GCED}} y \href{\linkPortalGIA}{\acrshort{GIA}}) requieren \textbf{240 créditos ECTS}.
\end{infoBox}

\subsection{Precio de los créditos}

El \textbf{precio por crédito aumenta} con cada matrícula en la misma asignatura.

\begin{curiosityBox}
    El precio por crédito de la \acrshort{UDC} es igual al del resto de universidades del Sistema Universitario de Galicia. Este se publica en el \acrfull{DOG} y puede ser consultado desde la \href{https://www.udc.es/es/matricula/}{página web de la UDC}.
\end{curiosityBox}

\subsection{Reconocimiento de créditos}

Es posible \textbf{obtener créditos} por actividades extracurriculares, cursos de idiomas, voluntariado o incluso actividades deportivas.

\begin{infoBox}
    Las actividades concretas que otorgan créditos pueden consultarse en la \href{https://www.udc.es/gl/normativa/academica/}{\textbf{normativa académica sobre reconocimiento de créditos}}.
\end{infoBox}

\subsection{Restricciones de créditos}

Los créditos también resultan de \textbf{gran importancia} para la realización de \textbf{actividades obligatorias} como las prácticas en empresas o el \acrfull{TFG}.

\subsection{Requisitos de permanencia}

Para poder \textbf{continuar los estudios} en la \acrshort{UDC} se requiere un \textbf{rendimiento mínimo}:

\begin{itemize}
    \item Alcanzar \textbf{12 créditos} \acrshort{ECTS} en el \textbf{primer año}.

    \item Alcanzar \textbf{60 créditos} \acrshort{ECTS} \textbf{cada tres cursos académicos}.
\end{itemize}

\begin{warningBox}
    En el caso de no cumplir con el mínimo de créditos tras el primer curso se puede solicitar la permanencia mediante una instancia al decano. En caso favorable, será obligatorio aprobar un mínimo de \textbf{36 créditos} \acrshort{ECTS} en el siguiente curso académico.
\end{warningBox}

\begin{infoBox}
    Las cantidades de créditos indicados corresponden a una matrícula a tiempo completo. En el caso de una matrícula a tiempo parcial correspondería la \textbf{mitad} de cada cantidad.
\end{infoBox}

\begin{importantBox}
    Todas las condiciones, procedimientos y excepciones están recogidas en la \href{https://www.udc.es/gl/normativa/academica/}{\textbf{normativa de gestión académica}}. Se recomienda encarecidamente su consulta.
\end{importantBox}


\section{Asignaturas}

\subsection{Guías docentes}

Las \href{https://guiadocente.udc.es/guia_docent/index.php?centre=614&ensenyament=614G01&consulta=assignatures&idioma=cast}{\textbf{guías docentes}} recogen la información sobre las asignaturas de los grados y másteres. Esta información incluye al profesorado, los contenidos, la metodología de evaluación y una breve descripción de la asignatura.  

\begin{curiosityBox}
    Las guías docentes de un curso son aprobadas a finales del curso académico anterior, por lo que pueden ser consultadas antes de realizar la matrícula.
\end{curiosityBox}

\subsection{Resultados}

Para cada asignatura de los últimos cursos anteriores es posible consultar el número total de \textbf{aprobados} (~\textcolor{resultPassed}{\fcolorbox{black}{resultPassed}{\rule{5pt}{5pt}}}~), \textbf{suspensos} (~\textcolor{resultFailed}{\fcolorbox{black}{resultFailed}{\rule{5pt}{5pt}}}~) y \textbf{no presentados} (~\textcolor{resultNotPresented}{\fcolorbox{black}{resultNotPresented}{\rule{5pt}{5pt}}}~) en el \href{\linkPortalEstudos}{portal de estudios}.

\FloatBarrier
\begin{figure}[htp]
    \centering
    \frame{\includegraphics[width=0.8\linewidth]{figures/asignaturas/resultados.png}}
\end{figure}
\FloatBarrier

\subsection{Tutorías} 

El profesorado cuenta con horarios dedicados a atender a los alumnos para resolver dudas que no se han resuelto en clase o para recibir atención más personalizada.  

Los \textbf{horarios y lugares de tutorías} de los profesores de la \acrshort{FIC} pueden ser consultados en la \href{https://www.udc.es/es/centros_departamentos_servizos/centros/titorias/?codigo=614}{página de tutorías de la \acrshort{UDC}} o en \href{https://espazos.udc.es/centers/614/tutorials}{Espazos \acrshort{UDC}}.

\begin{infoBox}
    Algunos profesores solicitan concertar una cita previa para sus tutorías, otros prefieren que, en caso de dudas puntuales, se les envíe dicha duda en un mensaje por Teams o correo electrónico. Es recomendable preguntar al profesorado por sus preferencias.     
\end{infoBox}
 
\subsection{Evaluación docente}

Al finalizar cada cuatrimestre, se podrán realizar unas \textbf{encuestas anónimas} para \textbf{valorar al profesorado} de cara a mejorar las clases y las asignaturas. Existe una encuesta por cada docente de las asignaturas en las que se está matriculado y debe llegarse a un \textbf{mínimo de encuestas} para que se tengan en cuenta. Estas encuestas tienen gran repercusión en la \textbf{reputación} del profesorado, por lo que es importante que se realicen de manera objetiva y sincera. 

\subsection{Evaluación por compensación}

Si a un estudiante solo le falta una asignatura para obtener el título y cumple unos determinados requisitos, puede solicitar su compensación y completar el estudio.

\begin{infoBox}
    No se pueden compensar ni las prácticas en empresa ni el \acrshort{TFG}.
\end{infoBox}
